\documentclass{article}
\usepackage[utf8]{inputenc}
\usepackage[T1]{fontenc}
\usepackage{geometry}
\usepackage{graphicx}
\usepackage{float}
\usepackage{hyperref}
\usepackage[english,indonesian]{babel}

\title{Optimasi Pemrosesan Template LaTeX Menggunakan Pendekatan Berbasis AI}
\author{Academic Support Agent}
\date{\today}

\begin{document}
\maketitle

\begin{abstract}

\end{abstract}

\chapter{PENDAHULUAN}


\chapter{TINJAUAN PUSTAKA}


% Original content (commented out):
% \documentclass[twocolumn,10pt]{article}
\usepackage[letterpaper,left=1.95cm,right=1.95cm,top=2.5cm,bottom=3.16cm]{geometry} % Márgenes y demás
\input{Styles_Eje.sty}  %%Llama el archivo de los estilos

%%%%%%% ENCABEZADO %%%%%%%%%%%%%%%%%%%%

{\title{\fontsize{16}{16}
\vspace*{-0.7cm}
\textbf{
Título del Artículo en Español (Times New Roman 16pt, con mayúsculas la primera letra de cada palabra (excepto conectores))
} \\[0.2cm]
\fontsize{12}{12}
\textbf{Título del Artículo en Inglés (negrita, 12 pt, con mayúsculas la primera letra de cada palabra (excepto conectores))}\\[0.2cm]}
}

\author[1,*]{\fontsize{10pt}{10pt}\selectfont \textbf{Javier Andrés Almeida M.Sc.}}
\author[2]{\textbf{Andrés A. Navarro}}
\author[1,3]{\textbf{Yeinzon Rodríguez}}


\affil[1]{
\fontsize{8pt}{8pt}\selectfont Departamento de ingeniería biomédica, Universidad ECCI, Bogotá - Colombia}
\fontsize{8}{8}
\affil[2]{Departamento de Ciencias Básicas, Universidad Santo Tomás, Bucaramanga, Colombia.}
\affil[3]{Centro de Investigaciones en Ciencias Básicas y Aplicadas, Universidad Antonio Nariño, Bogotá D.C., Colombia \vspace{0.2cm}} 
\affil[*]{jalmeidam@ecci.edu.co}



%%%%%%%%%%%%%%%%%%%%%%%%%%%%%%%%% INICIO DEL DOCUMENTO
%%%%%%%%%%%%%%%%%%%%%%%%%%%%%%%%%%%%%%%%%%%%%%%%%%%%%%%%%%%%%%%%%%%%%%%


\begin{document}
\input{Abstract}
\thispagestyle{alim}


\section{INTRODUCCIÓN}

\justifying
Este documento es un tutorial de cómo elaborar un artículo derivado de investigación y este archivo Word se constituye en una plantilla para presentar el artículo a la Revista EJE \cite{jacobs1996regulation}.

Para su selección, la revista EJE recibe artículos en idioma español o inglés, los artículos en español deben tener un abstract en inglés y viceversa.  Los artículos sometidos a consideración de la revista deben ser inéditos, es decir, que no se hayan publicado en otros medios, electrónicos o   impresos y que no hayan sido sometidos a otra revista.

Todo artículo de investigación debe tener las siguientes secciones: introducción, materiales y método, resultados, análisis de resultados, conclusión, agradecimientos y referencias. La \textbf{Introducción} es la parte donde el autor brinda una visión resumida y concisa del objetivo del artículo y de los trabajos previos e investigaciones que se han realizado sobre el tema. Esta sección debe incluir suficiente información de fondo, haga una exposición clara del problema, la literatura relevante sobre el tema y el enfoque, la aproximación o solución propuesta. Los objetivos del manuscrito deben ser claramente establecidos. La introducción no debe contener resultados o conclusiones. La introducción debe ser comprensible para colegas de una amplia gama de disciplinas científicas. 

\section{MATERIALES Y MÉTODO}

Esta sección explica cómo se hizo la investigación, debe ser lo suficientemente completa, proporcionar suficiente detalle como para permitir que el trabajo sea repetido y validado por otros. Sin embargo, sólo los procedimientos verdaderamente nuevos deben ser descritos con detalle, los procedimientos publicados anteriormente deben ser citados, y mencionar brevemente las modificaciones importantes de procedimientos ya publicados. Se debe utilizar subtítulos, ejemplo:

\subsection{Muestra}
\subsubsection{Ejemplo de sub sub capítulo}
\subsection{Descripción del Diseño del experimento }
\subsection{Toma de datos}
\subsection{Análisis de información y recursos utilizados}

Los métodos de uso general no necesitan ser descritos en detalle. Se debe evitar el uso de términos ambiguos, tales como frecuentemente, regularmente y periódicamente. No especificar marcas comerciales de productos ni modelos específicos de equipos. Se debe redactar en tiempo pasado y en tercera persona, por ejemplo: se realizó, se midió, se mezcló, etc.).

\section{RESULTADOS Y DISCUSIÓN}

Los resultados deben expresar lo obtenido de los experimentos o pruebas descritos en los materiales y métodos, se presentan en una secuencia lógica en el texto, utilizando tablas y figuras, se debe evitar la presentación repetida de los mismos datos en diferentes formas (al final de este documento se explica cómo presentar ecuaciones, figuras y tablas). Los resultados y discusión deben ser explicados y descritos en tiempo pasado \cite{lewismodeling}.

La discusión debe considerar los resultados en relación con las hipótesis formuladas en la introducción y la comparación del estudio con el contexto de otros trabajos. 

\section{CONCLUSIÓN}

Las conclusiones son obligatorias y deben ser claras. Su contenido no debería duplicar substancialmente el resumen. Deben expresar el balance final de la investigación o la aplicación del conocimiento o temática tratada. Se discute sobre las implicaciones del estudio y la relevancia que tiene para el área del conocimiento. Se sugiere no concluir más cosas de las que los resultados permitan.  En esta sección se suelen mencionar también los trabajos futuros que se pueden realizar en el tema con sus respectivas recomendaciones.

\section{AGRADECIMIENTOS}

Deben ser breves. Esta sección es obligatoria para artículos de investigación, en esta parte del artículo el autor hace un reconocimiento a las personas o instituciones que le ayudaron en sus investigaciones.  Se citan becas e instituciones que financian la investigación: firmas comerciales, entidades oficiales o privadas, asociaciones de profesionales y operarios.  Esta sección es opcional para artículos de reflexión. No es una sección numerada.

\section*{INFORMACIÓN ADICIONAL DEL FORMATO DE ARTÍCULO}

Para presentar los principales hallazgos del estudio, las tablas y los gráficos son claves, pues en ellos se ilustran claramente los resultados obtenidos.

En el cuerpo del artículo se utiliza letra tipo Times New Roman 10 puntos, excepto en el título. Antes de cada título se deja doble espacio y después de cada título, se dejará un espacio sencillo antes de iniciar el párrafo. La extensión de un artículo depende del tipo de artículo (científico, de revisión o corto) sin embargo, no será mayor a veinte (20) páginas a doble columna y espacio sencillo, las páginas deben ser de tamaño carta.

\section*{ECUACIONES}

Se deben numerar consecutivamente en paréntesis normal, en el margen derecho, sin hacer referencia a la sección donde se encuentra, es decir, debe haber una numeración independiente para las ecuaciones.  Para escribir la ecuación, utilice el editor de ecuaciones.  Es importante que los símbolos se definan antes o inmediatamente después de que aparece la ecuación.

Unidades, siempre debe haber un espacio después de todas las unidades (excepto para temperatura y porcentajes) y deben estar sin formato (no cursiva ni negrilla).

Ejemplo:

\begin{equation}
V(t) = V_{1}*e^{-\infty t } - \sum_{k-1}^{100}M_{k}*e^{-t}
\end{equation}

\section*{Figuras y Tablas}

Todas deben llevar numeración arábiga de acuerdo con su orden de aparición y, al igual que las ecuaciones, se hará referencia a ellas en el texto más cercano a la misma. Deberán ser nítidas, las fotografías y figuras deben ser originales, a color, en blanco y negro o tonos de gris, de acuerdo a la preferencia del autor, con una resolución mínima de 200dpi (puntos por pulgada); el título debe ser escrito en forma clara y que explique su contenido, ubicado en la parte inferior, justificado a la izquierda, en letra tipo Times New Roman a 10 pts., como se observa en la Fig.\ref{fig:Ejemplo_Fig}. En las tablas, el título se presentará en la parte superior, con alineación justificada, en letra tipo Times New Roman a 9 pts, como se muestra en Tabla \ref{Tab:Ejemplo_Tabla}.

Las tablas son flotantes, por lo que su referencia dentro del texto no debe depender de su ubicación para fines de organización. Sin embargo, en lo posible las tablas deben ir en el texto siguiente al que hace referencia a la tabla. Deben estar distribuidas a lo largo del texto de manera que no haya ninguna después de la sección de Referencias \cite{allen1991chaos}.

Ejemplo figura:

\begin{figure}[H]
    \centering
    \includegraphics[scale=1]{Figuras/Ejemplo_1.png}
    \caption{La leyenda debe explicar claramente la figura correspondiente}
    \label{fig:Ejemplo_Fig}
\end{figure}


\begin{table}[H]
\caption{Tipos de letra, justificación y tamaño}
\begin{tabular}{llll}
\hline
Letra    & Tamaño & Letra & Justificación \\ \hline
Titulo   & 16pt   & Arial & Centrada      \\ \hline
Autores  & 10pt   & Arial & Centrada      \\ \hline
Texto    & 10pt   & Arial & Justificada   \\ \hline
Ecuación & 10pt   & Arial & Justificada   \\ \hline
Figura   & 10pt   & Arial & Justificada   \\ \hline
Tabla    & 10pt   & Arial & Justificada   \\ \hline
\end{tabular}
\label{Tab:Ejemplo_Tabla}
\end{table}









\bibliographystyle{ieeetr}
\bibliography{Bibliografia}

\textcolor{red}{Seguir las normas IEEE. Deben insertarse dentro del cuerpo del artículo, con números arábigos entre corchetes [1], consecutivamente, en la medida en que van siendo mencionadas. No es una sección numerada.}

\textcolor{red}{
En la lista, al final, se organizan consecutivamente en el orden de citación (nunca en orden alfabético).  No deben incluirse referencias que no estén citadas en el artículo.  Igualmente, no deben citarse trabajos en revisión, que no hayan sido publicados.}

%Descomente las últimas líneas al final







































































































\end{document}


\end{document}